\documentclass[11pt,a4paper]{article}
\usepackage[margin=2.55cm]{geometry}
\usepackage{amsmath,amssymb,amsthm}
\usepackage{booktabs}
\usepackage[hidelinks]{hyperref}
\hypersetup{
  pdftitle={Positive Spectral Measures as Operational Resource Boundaries},
  pdfauthor={Lluis Eriksson},
  pdfsubject={Correlations, resolvents, mass gaps, and rate inheritance}
}

\newtheorem{theorem}{Theorem}[section]
\newtheorem{proposition}[theorem]{Proposition}
\newtheorem{corollary}[theorem]{Corollary}
\newtheorem{lemma}[theorem]{Lemma}
\theoremstyle{definition}
\newtheorem{definition}[theorem]{Definition}
\theoremstyle{remark}
\newtheorem{remark}[theorem]{Remark}

\newcommand{\C}{\mathbb C}
\newcommand{\R}{\mathbb R}
\newcommand{\N}{\mathbb N}
\newcommand{\supp}{\operatorname{supp}}
\newcommand{\spec}{\operatorname{spec}}
\newcommand{\Span}{\operatorname{span}}
\newcommand{\kB}{k_{\mathrm B}}

\title{Positive Spectral Measures as Operational Resource Boundaries:\\
Correlations, Resolvents, Mass Gaps, and Rate Inheritance}
\author{Lluis Eriksson\\Independent Researcher\\
\href{mailto:lluiseriksson@gmail.com}{lluiseriksson@gmail.com}}
\date{August 2026}

\begin{document}
\maketitle

\begin{abstract}
We give a common positive-measure formulation of four statements that are
often discussed separately: exponential Euclidean clustering, a Hamiltonian or
transfer-operator mass gap, Stieltjes resolvent positivity, and the exponential
suppression of a dissipative rate transmitted through a buffer.  For a
positive self-adjoint Hamiltonian and a vacuum-orthogonal vector, we prove that
absence of spectral weight below $m$ is equivalent to an exponential
correlation envelope with exponent $m$.  Uniform support exclusion on a total
family is then equivalent to an operatorial mass gap.  In discrete transfer
time, the same statement becomes a Hausdorff support problem: Hankel and
localizing moment matrices provide finite falsifiers of an overstated gap.  We
separate this closed spectral theorem from the genuinely model-dependent
open-system interface.  Once a rate is bounded by the square of a mediated
amplitude, its exponent doubles.  Conversely, a two-sided rate--correlator
comparison on a total positive family recovers a transfer gap.  We also correct
an operational ambiguity: a lower bound on maintenance power defines an
impossibility horizon, not an affordability horizon; a sharp resource boundary
requires a constructive upper bound.  Reflection-positive lattice gauge theory
and prime-built resolvent programmes share this spectral-measure architecture,
but no implication between the Yang--Mills mass-gap problem and the Riemann
hypothesis follows.  Exact rational certificates, an independent checker, and
deterministic numerical stress tests accompany the paper.
\end{abstract}

\section{Introduction}

A spatial buffer can suppress the amplitude by which a disturbance reaches a
protected system.  If the dissipative rate is quadratic in that amplitude, an
amplitude law $e^{-m\ell}$ becomes a rate law $e^{-2m\ell}$.  This mechanism is
elementary in exactly solvable band-gap models, yet its interpretation touches
three harder questions: what spectral statement is actually being measured,
when a family of such measurements certifies a mass gap, and when a suppressed
rate makes coherence maintenance affordable.

The motivating calculation is the quasi-free gapped-buffer model of
Ref.~\cite{Eriksson0064}, where an exact Liouvillian computation gives a
frequency-resolved rate exponent twice the evanescent-amplitude exponent.  A
companion stress test~\cite{Eriksson0070} shows why ceiling and floor quantities
must not be interchanged.  Here we extract the model-independent mathematical
core and identify the interfaces that remain physical input.

The organizing object is not a decay curve but a positive spectral measure.
Let $H\geq0$ be a self-adjoint Hamiltonian, let $\Omega$ be a vacuum, and let
$\psi\perp\Omega$.  The measure
\[
  \mu_\psi(B)=\langle\psi,E_H(B)\psi\rangle
\]
simultaneously generates the Euclidean correlation
\[
  C_\psi(t)=\int e^{-t\lambda}\,d\mu_\psi(\lambda)
\]
and the one-point resolvent
\[
  S_\psi(x)=\int\frac{d\mu_\psi(\lambda)}{\lambda+x}.
\]
The first is a Laplace transform, the second a Stieltjes transform, and a mass
gap is a lower support edge.  This triangle is the precise common structure
behind the correlations used in constructive field theory, the transfer
operators used in reflection-positive reconstruction, and the positive
resolvents used in spectral criteria.

The contributions are as follows.
\begin{enumerate}
  \item We prove a vectorwise equivalence between a spectral support edge and
  an exponential correlation exponent, and lift it to an operator gap through
  a total-family criterion.
  \item We formulate the discrete transfer version as a Hausdorff moment
  problem and give finite localizing-matrix falsifiers for claimed support
  edges.
  \item We isolate a rate-mediation lemma that explains exponent doubling, and
  a converse rate-to-gap criterion whose strong hypotheses are fully visible.
  \item We replace a single ``resource boundary'' by two operational horizons:
  certified impossibility and constructive affordability.
  \item We state exactly how the framework can be consumed by
  Osterwalder--Schrader reconstruction and lattice Yang--Mills, and exactly why
  the analogous Stieltjes architecture in prime-resolvent programmes does not
  imply a Yang--Mills--Riemann equivalence.
\end{enumerate}

The spectral results are applications of standard functional calculus.  The
novelty claimed here is the operational assembly, the two-sided audit of what a
rate can certify, and a reproducible support-certificate layer---not a new proof
of the spectral theorem, the Yang--Mills mass gap, or the Riemann hypothesis.

\section{The positive spectral triangle}
\label{sec:triangle}

Let $\mathcal H$ be a complex Hilbert space and let $H$ be a nonnegative
self-adjoint operator.  We assume a unique normalized vacuum,
\[
  \ker H=\C\Omega,\qquad \|\Omega\|=1.
\]
For $\psi\in\Omega^\perp$, let $\mu_\psi$ be its finite spectral measure,
\[
  \mu_\psi(B)=\langle\psi,E_H(B)\psi\rangle,
  \qquad \mu_\psi([0,\infty))=\|\psi\|^2.
\]

\begin{proposition}[Laplace and Stieltjes faces]
\label{prop:faces}
For $t>0$ and $x>0$, define
\begin{align}
  C_\psi(t)&=\langle\psi,e^{-tH}\psi\rangle,
  &S_\psi(x)&=\langle\psi,(H+x)^{-1}\psi\rangle.
\end{align}
Then
\begin{align}
  C_\psi(t)&=\int_{[0,\infty)}e^{-t\lambda}\,d\mu_\psi(\lambda),
  \label{eq:laplace-face}\\
  S_\psi(x)&=\int_{[0,\infty)}\frac{d\mu_\psi(\lambda)}{\lambda+x}.
  \label{eq:stieltjes-face}
\end{align}
Moreover, $C_\psi$ is completely monotone on $(0,\infty)$ and $S_\psi$ is a
Stieltjes function.  We use the continuous extension
$C_\psi(0)=\|\psi\|^2$ whenever $t=0$ appears below.
\end{proposition}

\begin{proof}
Equations~\eqref{eq:laplace-face} and \eqref{eq:stieltjes-face} are the
spectral functional calculus.  For every $n\in\N$ and $t>0$,
\[
  (-1)^n C_\psi^{(n)}(t)
  =\int\lambda^n e^{-t\lambda}\,d\mu_\psi(\lambda)\geq0.
\]
The Stieltjes representation is Eq.~\eqref{eq:stieltjes-face} itself.
\end{proof}

The next theorem identifies the exponential exponent without assuming a
discrete spectrum.

\begin{theorem}[Support--decay equivalence]
\label{thm:support-decay}
Let $m\geq0$.  The following are equivalent:
\begin{enumerate}
  \item $\mu_\psi([0,m))=0$;
  \item $C_\psi(t)\leq\|\psi\|^2e^{-mt}$ for every $t\geq0$;
  \item there exists a finite $K_\psi$ such that
  $C_\psi(t)\leq K_\psi e^{-mt}$ for every $t\geq0$.
\end{enumerate}
\end{theorem}

\begin{proof}
If the measure is supported in $[m,\infty)$, then
$e^{-t\lambda}\leq e^{-mt}$ and (2) follows by integration.  Statement (2)
implies (3).  Conversely, suppose (3) holds but
$\mu_\psi([0,m))>0$.  Since
$[0,m)=\bigcup_{n\geq1}[0,m-1/n]$ after discarding empty terms, there is a
$\delta>0$ with $a:=\mu_\psi([0,m-\delta])>0$.  Positivity gives
\[
  C_\psi(t)\geq a e^{-(m-\delta)t}.
\]
The ratio of this lower bound to $K_\psi e^{-mt}$ is
$(a/K_\psi)e^{\delta t}$, which exceeds one for sufficiently large $t$, a
contradiction.
\end{proof}

\begin{corollary}[Resolvent bound from a support edge]
\label{cor:resolvent-bound}
Under the equivalent conditions of Theorem~\ref{thm:support-decay},
\[
  S_\psi(x)\leq\frac{\|\psi\|^2}{m+x},\qquad x>0.
  \label{eq:resolvent-bound}
\]
\end{corollary}

\begin{proof}
On the support of the measure,
$(\lambda+x)^{-1}\leq(m+x)^{-1}$.
\end{proof}

A single vector can be blind to a gapless sector.  The correct upgrade is a
total-family condition.

\begin{theorem}[Total-family gap criterion]
\label{thm:total-family}
Let $\mathcal D\subset\Omega^\perp$ have dense linear span in
$\Omega^\perp$.  If there is an $m>0$ such that every $\psi\in\mathcal D$
satisfies any of the equivalent conditions of
Theorem~\ref{thm:support-decay}, then
\[
  H|_{\Omega^\perp}\geq m
\]
in the spectral, equivalently quadratic-form, sense.  Conversely, this operator
gap implies the support condition for every $\psi\perp\Omega$.
\end{theorem}

\begin{proof}
Let $P=E_H([0,m))$.  The hypothesis says $P\psi=0$ for every
$\psi\in\mathcal D$.  As $P$ is bounded, it vanishes on the closure of
$\Span\mathcal D=\Omega^\perp$.  Hence the spectrum of $H$ on
$\Omega^\perp$ lies in $[m,\infty)$.  The converse follows immediately from
spectral support.
\end{proof}

\begin{remark}[Uniform exponent versus uniform prefactor]
Theorem~\ref{thm:total-family} needs a common exponent $m$, but it does not
need a common prefactor $K_\psi$.  Volume-uniform physical applications are
harder because one must also control how the family, normalization, and
finite-volume reconstruction behave as the volume changes.
\end{remark}

\begin{corollary}[Off-diagonal transfer amplitude]
\label{cor:offdiagonal}
Under the operator gap of Theorem~\ref{thm:total-family}, let
$P_\perp=I-|\Omega\rangle\langle\Omega|$.  Then
\[
  \|e^{-tH}P_\perp\|\leq e^{-mt},
\]
and for all $\phi,\psi\perp\Omega$,
\[
  |\langle\phi,e^{-tH}\psi\rangle|
  \leq e^{-mt}\|\phi\|\,\|\psi\|.
  \label{eq:offdiagonal}
\]
\end{corollary}

\begin{proof}
The operator-norm bound is the spectral calculus on $\Omega^\perp$; the
matrix-element bound follows from Cauchy--Schwarz.
\end{proof}

\section{Transfer operators and finite support certificates}
\label{sec:transfer}

Let $T$ be a positive self-adjoint contraction with $T\Omega=\Omega$.  For
$\psi\perp\Omega$, write $\nu_\psi$ for the spectral measure of $T$ on
$[0,1]$ and
\[
  b_n(\psi)=\langle\psi,T^n\psi\rangle
  =\int_{[0,1]}r^n\,d\nu_\psi(r).
  \label{eq:transfer-moments}
\]
Thus the transfer correlation sequence is a Hausdorff moment sequence.  The
continuous and discrete pictures agree under $T=e^{-aH}$: a support edge
$H\geq m$ becomes $T\leq\rho=e^{-am}$ on $\Omega^\perp$.

\begin{proposition}[Discrete support--decay equivalence]
\label{prop:discrete-support}
For $0<\rho<1$, the following are equivalent:
\begin{enumerate}
  \item $\nu_\psi((\rho,1])=0$;
  \item $b_n(\psi)\leq\|\psi\|^2\rho^n$ for every $n\in\N$;
  \item $b_n(\psi)\leq K_\psi\rho^n$ for every $n\in\N$ and some finite
  $K_\psi$.
\end{enumerate}
\end{proposition}

\begin{proof}
The forward implication is immediate.  If the measure has positive mass on
$[\rho+\delta,1]$ for some $\delta>0$, that mass gives a lower bound
$a(\rho+\delta)^n$, contradicting any finite multiple of $\rho^n$ as
$n\to\infty$.
\end{proof}

Finite moments cannot prove an infinite support theorem, but they can refute a
false support claim exactly.  For $N\geq0$ define
\begin{align}
  H_N&=(b_{i+j})_{0\leq i,j\leq N},\\
  L_N^\theta&=(\theta b_{i+j}-b_{i+j+1})_{0\leq i,j\leq N}.
  \label{eq:moment-matrices}
\end{align}

\begin{proposition}[Finite support falsifier]
\label{prop:finite-falsifier}
If $\nu_\psi$ is supported in $[0,\theta]$, then
$H_N\succeq0$ and $L_N^\theta\succeq0$ for every $N$.  Consequently, a
single vector $q$ with
\[
  q^{\mathsf T}L_N^\theta q<0
\]
is an exact finite falsifier of the claimed support edge $\theta$.
\end{proposition}

\begin{proof}
For $p_q(r)=\sum_{i=0}^Nq_ir^i$,
\begin{align}
  q^{\mathsf T}H_Nq
  &=\int p_q(r)^2\,d\nu_\psi(r)\geq0,\\
  q^{\mathsf T}L_N^\theta q
  &=\int(\theta-r)p_q(r)^2\,d\nu_\psi(r)\geq0.
\end{align}
\end{proof}

Passing all finite tests is only consistency.  A proof of support requires the
whole hierarchy plus the measure/reconstruction hypotheses, or an independent
analytic argument.  This one-sided epistemology is valuable: negative
certificates are decisive, positive finite certificates are not terminal.

\section{Rate inheritance and its converse}
\label{sec:rate}

The spectral triangle is closed.  A dissipative statement requires an
additional interface connecting the system, the buffer, and the sink.  Let
$M_{LR}(\ell,\omega)$ be a boundary-to-boundary amplitude at system Bohr
frequency $\omega$, and let $\kappa_\ell(\omega)$ be the induced rate.

\begin{definition}[Mediation interface]
We say that a model has a $(\Gamma,r)$ mediation upper interface when
\[
  \kappa_\ell(\omega)
  \leq\Gamma(\omega)|M_{LR}(\ell,\omega)|^2
  +r_T(\ell,\omega),
  \label{eq:mediation-interface}
\]
where $r_T$ contains the thermal floor and controlled open-system reduction
errors.
\end{definition}

\begin{lemma}[Exponent squaring]
\label{lem:exponent-squaring}
If Eq.~\eqref{eq:mediation-interface} holds and
\[
  |M_{LR}(\ell,\omega)|
  \leq C(\omega)(1+\ell)^p e^{-m(\omega)\ell},
  \label{eq:amplitude-bound}
\]
then
\[
  \kappa_\ell(\omega)
  \leq\Gamma(\omega)C(\omega)^2(1+\ell)^{2p}
  e^{-2m(\omega)\ell}+r_T(\ell,\omega).
  \label{eq:rate-bound}
\]
\end{lemma}

\begin{proof}
Square Eq.~\eqref{eq:amplitude-bound} and substitute it into
Eq.~\eqref{eq:mediation-interface}.
\end{proof}

The lemma is intentionally simple.  It identifies the actual interacting
frontier: proving a volume-uniform amplitude bound and a width-uniform
open-system reduction in the presence of the sink.  Static exponential
clustering alone does not automatically control the real-time spectral density
or the state driven by that sink.

There is a converse, but only after excluding blindness and cancellations.

\begin{theorem}[Operational rate-to-gap criterion]
\label{thm:rate-to-gap}
Let $T$ be as in Section~\ref{sec:transfer}, and let
$\mathcal D\subset\Omega^\perp$ be total.  Suppose that for every
$\psi\in\mathcal D$ the measured channel obeys, for all $n$,
\[
  \gamma_- b_n(\psi)^2\leq\kappa_n(\psi)
  \leq\gamma_+ b_n(\psi)^2,
  \qquad 0<\gamma_-\leq\gamma_+<\infty,
  \label{eq:two-sided-rate}
\]
and that a common $0<\rho<1$ satisfies
\[
  \kappa_n(\psi)\leq K_\psi\rho^{2n}
\]
with finite $K_\psi$.  Then $T|_{\Omega^\perp}\leq\rho$.  If
$T=e^{-aH}$, then $H|_{\Omega^\perp}\geq-a^{-1}\log\rho$.
\end{theorem}

\begin{proof}
The lower half of Eq.~\eqref{eq:two-sided-rate} gives
\[
  b_n(\psi)\leq\sqrt{K_\psi/\gamma_-}\,\rho^n.
\]
Proposition~\ref{prop:discrete-support} excludes transfer spectral weight above
$\rho$ for every $\psi\in\mathcal D$.  Totality and the projection argument of
Theorem~\ref{thm:total-family} give the operator bound.  Spectral calculus for
$T=e^{-aH}$ gives the Hamiltonian bound.
\end{proof}

\begin{remark}[Why the hypotheses are strong]
An upper rate bound alone cannot yield a gap: the probe may decouple from a
gapless sector, or interference may suppress a complex amplitude.  Reflection
positivity supplies nonnegative transfer correlations, while totality and the
lower comparison in Eq.~\eqref{eq:two-sided-rate} eliminate the two loopholes.
These are obligations, not consequences of observing one exponential curve.
\end{remark}

\section{A resource boundary has two horizons}
\label{sec:horizons}

Let $P_{\min}(\ell)$ be the minimum power required to maintain a target
coherence against the induced channel, within a specified operational model.
Thermodynamic arguments commonly provide a lower bound such as
$P_{\min}\geq\kB T\dot C_{\rm loss}$.  A small lower bound is not an upper
bound and therefore does not exhibit an affordable protocol.

\begin{definition}[Operational horizons]
For a power budget $P_{\max}$:
\begin{enumerate}
  \item the \emph{impossibility region} is where a proved lower bound on
  $P_{\min}$ exceeds $P_{\max}$;
  \item the \emph{affordability region} is where a constructive protocol gives
  an upper bound on $P_{\min}$ below $P_{\max}$.
\end{enumerate}
The interval between the two certified regions is an epistemic, not
necessarily physical, uncertainty region.
\end{definition}

\begin{proposition}[Two-horizon bracket]
\label{prop:horizon-bracket}
Assume an exactly exponential lane
$\kappa_\ell=\kappa_0e^{-\alpha\ell}$ and a two-sided power comparison
\[
  a_-\kappa_\ell\leq P_{\min}(\ell)\leq a_+\kappa_\ell,
  \qquad 0<a_-\leq a_+.
  \label{eq:power-comparison}
\]
When $a_-\kappa_0>P_{\max}$, maintenance is certified impossible for
\[
  \ell<\ell_{\rm imp}
  :=\frac1\alpha\log\frac{a_-\kappa_0}{P_{\max}}.
\]
When $a_+\kappa_0>P_{\max}$, the constructive protocol is certified affordable
for
\[
  \ell\geq\ell_{\rm aff}
  :=\frac1\alpha\log\frac{a_+\kappa_0}{P_{\max}}.
\]
The unresolved width is
\[
  \ell_{\rm aff}-\ell_{\rm imp}
  =\frac1\alpha\log\frac{a_+}{a_-}.
  \label{eq:horizon-width}
\]
\end{proposition}

\begin{proof}
Substitute the exponential rate into the lower and upper halves of
Eq.~\eqref{eq:power-comparison} and solve the two budget inequalities.
\end{proof}

Thermal activation or bypass channels replace the asymptotic zero by a floor.
The horizons then saturate or disappear; adding buffer no longer buys
arbitrarily cheap isolation.  This is a physical limitation, not a proof
failure.

\section{Reflection positivity and an operational mass-gap dictionary}
\label{sec:os}

Reflection positivity can turn a Euclidean measure into a Hilbert space, a
vacuum, and a positive transfer semigroup through the
Osterwalder--Schrader/GNS quotient~\cite{OS1973,OS1975,FILS1978}.  In that
setting a centered positive-time observable $A$ produces
$\psi_A=(A-\langle A\rangle)\Omega$, and its Euclidean two-point function is
the Laplace face of $\mu_A$.

\begin{corollary}[Conditional operational consequence of a mass gap]
\label{cor:ym-operational}
Suppose a reflection-positive lattice theory reconstructs a Hamiltonian with
gap $m>0$ on the vacuum-orthogonal sector.  For any gauge-invariant probe whose
boundary mediation amplitude is bounded by the corresponding transfer
amplitude, and whose open-system reduction satisfies
Eq.~\eqref{eq:mediation-interface}, the induced rate has exponential part at
most
\[
  \mathrm{poly}(\ell)e^{-2m\ell}
\]
up to lattice-spacing conventions, the thermal floor, and the reduction
remainder.
\end{corollary}

\begin{proof}
The gap gives the boundary-to-boundary amplitude exponent through
Corollary~\ref{cor:offdiagonal}, provided the centered boundary vectors lie in
$\Omega^\perp$.  Lemma~\ref{lem:exponent-squaring} gives the rate exponent.
\end{proof}

This corollary is a dictionary, not a Yang--Mills mass-gap proof.  In the
four-dimensional problem one must still construct the continuum theory,
establish reflection positivity and reconstruction in the required limit, prove a
volume-uniform gap for a total gauge-invariant family, and justify the
open-system mediation interface.  The current Eriksson Yang--Mills programme
contains a machine-checked strong-coupling lattice core, cluster-expansion
clustering, and explicit conditional interfaces, while keeping the continuum
and Clay boundary open~\cite{ErikssonProgramme}.  Rooted-tree and polymer
machinery may help control a two-boundary interacting amplitude, but abstract
summation does not produce the model-specific activity estimate.

The reverse direction is nevertheless conceptually useful.  If a
reflection-positive theory supplies a total positive family and a two-sided
rate identity of the form \eqref{eq:two-sided-rate}, uniform rate suppression
becomes a gap certificate through Theorem~\ref{thm:rate-to-gap}.  The theorem
states what would be sufficient; it does not assert that Yang--Mills currently
satisfies those hypotheses.

\section{Prime resolvents: the shared structure and its limit}
\label{sec:riemann}

The Stieltjes face of Proposition~\ref{prop:faces} also appears in spectral
approaches to the Riemann hypothesis.  With $\xi$ the completed zeta function,
pairing conjugate zeros in its Hadamard product gives under RH
\cite{TitchmarshZeta}
\[
  \mathcal S_\Xi(x)
  :=\frac{1}{2\sqrt{x}}\frac{\xi'}{\xi}
      \left(\frac12+\sqrt{x}\right)
  =\sum_{\gamma>0}\frac{1}{\gamma^2+x},
  \qquad x>\frac14,
  \label{eq:xi-stieltjes}
\]
with zeros counted by multiplicity.  The prime-resolvent
programme~\cite{PrimeResolvent} asks whether prime-built finite operators can
converge to that positive resolvent with explicit error budgets.

The common architecture is
\[
  \text{positive measure}
  \longrightarrow
  \text{Stieltjes resolvent}
  \longrightarrow
  \text{moment/localizing positivity}.
\]
It supports a common audit method: a negative finite Hankel or localizing
certificate can refute a proposed positive-measure or support representation,
whereas finitely many positive certificates establish consistency only.

No physical or logical equivalence follows.  In Yang--Mills, the measure is a
spectral measure of a reconstructed Hamiltonian and the gap concerns its lower
support on $\Omega^\perp$.  In the Riemann programme, positivity is tied to the
location of zeta zeros and the unresolved step is the convergence/continuation
of a prime-built resolvent.  Neither problem supplies the missing producer for
the other.  The statement that ``the universe constructed the primes'' is not
a mathematical consequence: primes are arithmetically defined.  What is
mathematically meaningful is that both problems can be tested through the same
positive-transform language.

\section{Exact finite certificates and reproducibility}
\label{sec:verification}

The repository contains a deterministic generator and an independent checker
at
\begin{center}
\href{https://github.com/lluiseriksson/aqft-split-inclusion-series/tree/agent/positive-spectral-resource-boundaries/verification/spectral_resource_boundary}
{\texttt{verification/spectral\_resource\_boundary/}}.
\end{center}
The exact layer uses rational arithmetic.  Consider the transfer atoms
\[
  r=(1/2,1/3,1/5),\qquad w=(2/7,3/7,2/7).
\]
For the correct edge $\theta=1/2$, the localizing matrix has the explicit
atomic Gram decomposition
\[
  L_2^{1/2}=\sum_k w_k(1/2-r_k)
  (1,r_k,r_k^2)^{\mathsf T}(1,r_k,r_k^2)\succeq0.
\]
For the false edge $\theta=2/5$, the polynomial
$p(r)=1-8r+15r^2$ vanishes at $1/3$ and $1/5$ but not at $1/2$.  The exact
witness is
\[
  (1,-8,15)L_2^{2/5}(1,-8,15)^{\mathsf T}=-\frac9{560}<0.
\]

The continuous model has excited energies $(2,3,5)$ with the same weights.
At $x=1$ its exact resolvent is
\[
  S(1)=\frac14\leq\frac13=\frac{\|\psi\|^2}{m+1},
  \qquad m=2.
\]
The tracked stress test also demonstrates that the overstated gap $m=2.1$
fails: at $t=30$ the ratio to its claimed envelope is approximately $5.739$.
Squaring the correlation and fitting $t=20,\ldots,40$ gives the log-rate slope
$-4.00000000012$, against the analytic value $-4$, with absolute error
$1.20\times10^{-10}$.

For the horizon example, $\alpha=4$, $P_{\max}=10^{-6}$,
$a_-=1$, and $a_+=4$ give
\[
  \ell_{\rm imp}=3.45388,\qquad
  \ell_{\rm aff}=3.80045,
\]
with uncertainty width $\log 4/4=0.346574$.  Table~\ref{tab:checks}
summarizes the audit.

\begin{table}[ht]
\centering
\caption{Deterministic verification outcomes.  ``Exact'' means rational
arithmetic with no floating-point decisions.}
\label{tab:checks}
\begin{tabular}{@{}lll@{}}
\toprule
Check & Arithmetic & Outcome \\
\midrule
$H_2$ moment positivity & exact & atomic Gram certificate \\
$L_2^{1/2}$ support test & exact & PSD Gram certificate \\
$L_2^{2/5}$ false edge & exact & witness $-9/560$ \\
$S(1)\leq(m+1)^{-1}$ & exact & $1/4\leq1/3$ \\
False gap $2.1$ & deterministic float & violated at $t=30$ \\
Squared-rate slope & deterministic float & error $1.20\times10^{-10}$ \\
Two-horizon ordering & deterministic float & $\ell_{\rm imp}<\ell_{\rm aff}$ \\
\bottomrule
\end{tabular}
\end{table}

Run from the repository root:
\begin{verbatim}
python verification/spectral_resource_boundary/generate_artifacts.py --check
python verification/spectral_resource_boundary/verify_artifact.py
\end{verbatim}
The JSON artifact is byte-deterministic.  Its scope is finite certification and
toy-model falsification; it is not numerical evidence for a continuum mass gap
or for RH.

\section{Open producers and a research programme}
\label{sec:open}

The framework separates reusable consumers from model-specific producers.
This makes the next steps testable.

\paragraph{Interacting buffers.}
Prove a two-boundary amplitude estimate in the state actually perturbed by the
sink, not only static ground-state clustering.  A viable first target is a
finite-range commuting chain at low temperature with a uniform mixing or
stability estimate.  Lieb--Robinson bounds~\cite{LiebRobinson,Poulin2010} and
gap--clustering results~\cite{HastingsKoma} control parts of the problem; the
uniform Davies limit and sink back-reaction remain separate obligations.

\paragraph{Polymer route.}
For a weakly interacting or strong-coupling regime, introduce a two-marked
activity whose clusters connect the protected and dissipative boundaries.  A
source estimate
\[
  \sum_{\Gamma:L\leftrightarrow R}|w(\Gamma)|
  \leq C e^{-m\ell}
\]
would feed the abstract rooted-tree and Koteck\'y--Preiss summation machinery.
The hard step is the physical activity estimate, not closing the already
abstracted geometric series.

\paragraph{Reflection-positive gauge models.}
First prove the forward statement of Corollary~\ref{cor:ym-operational} in an
exactly controlled two-dimensional or finite transfer model.  Only then test
which ingredients survive volume and continuum limits.  An operational
reverse theorem requires a total gauge-invariant probe family and a uniform
lower comparison, not a selected Wilson loop.

\paragraph{Moment falsifiers.}
Use exact or interval localizing matrices to challenge claimed transfer edges
and resolvent support windows.  A negative certificate is a genuine finite
obstruction.  A long sequence of positive certificates is a useful stress
test, not a substitute for the infinite theorem.

\section{Discussion}

The framework clarifies four distinctions.

First, a gap is not merely a fitted slope.  It is a support exclusion for a
positive spectral measure.  The fitted slope becomes decisive only when the
observable family is total and the correlation or rate is tied to that measure
without blind sectors.

Second, rate inheritance is not identical to clustering.  The spectral theorem
controls Euclidean correlations; an induced dissipative rate additionally
requires a dynamical mediation identity and control of the sink.  Once that
interface is present, the doubled exponent is unavoidable and elementary.

Third, operational classicality has two certification directions.  A power
lower bound can certify impossibility.  Affordability needs an upper bound or a
protocol.  Collapsing the two produces a boundary sharper than the evidence.

Fourth, common mathematics is not common ontology.  Yang--Mills transfer
spectra and prime-resolvent criteria both use positive measures and Stieltjes
transforms.  This licenses shared lemmas, certificates, and falsification
protocols; it does not license a physical identification of primes with field
excitations or a logical implication between open problems.

\section{Conclusion}

Positive spectral measures provide a compact, auditable language for
correlation decay, transfer gaps, resolvent positivity, and mediated rates.  A
support edge is equivalent to a vectorwise exponential exponent; a common edge
on a total family is an operator mass gap.  Finite moment matrices can refute
false support claims.  A squared mediation law doubles the exponent, while a
two-sided rate identity can, under positivity and totality, recover the gap.
Finally, a resource boundary is generally an interval between impossibility
and affordability horizons, not a single point read from a lower bound.

The result is a theorem-and-interface architecture.  It advances the
operational programme by closing the universal spectral layer and making every
remaining physical producer explicit.  That is also the appropriate level at
which the Yang--Mills and prime-resolvent programmes meet: shared positive
spectral technology, with their hardest claims still separate and open.

\section*{Acknowledgements}
The exact certificate scripts and all numerical values in
Section~\ref{sec:verification} are included with the source.  Computational
assistance was used for drafting and verification; all claim boundaries and
remaining errors are the author's responsibility.

\begin{thebibliography}{99}

\bibitem{Bernstein1929}
S.~Bernstein, \emph{Sur les fonctions absolument monotones}, Acta Math.
\textbf{52}, 1--66 (1929).

\bibitem{Widder1941}
D.~V.~Widder, \emph{The Laplace Transform}, Princeton University Press (1941).

\bibitem{ReedSimon}
M.~Reed and B.~Simon, \emph{Methods of Modern Mathematical Physics I:
Functional Analysis}, revised edition, Academic Press (1980).

\bibitem{CombesThomas}
J.~M.~Combes and L.~Thomas, \emph{Asymptotic behaviour of eigenfunctions for
multiparticle Schr\"odinger operators}, Commun. Math. Phys. \textbf{34},
251--270 (1973).

\bibitem{LiebRobinson}
E.~H.~Lieb and D.~W.~Robinson, \emph{The finite group velocity of quantum spin
systems}, Commun. Math. Phys. \textbf{28}, 251--257 (1972).

\bibitem{HastingsKoma}
M.~B.~Hastings and T.~Koma, \emph{Spectral gap and exponential decay of
correlations}, Commun. Math. Phys. \textbf{265}, 781--804 (2006).

\bibitem{Poulin2010}
D.~Poulin, \emph{Lieb--Robinson bound and locality for general Markovian
quantum dynamics}, Phys. Rev. Lett. \textbf{104}, 190401 (2010).

\bibitem{Davies1974}
E.~B.~Davies, \emph{Markovian master equations}, Commun. Math. Phys.
\textbf{39}, 91--110 (1974).

\bibitem{OS1973}
K.~Osterwalder and R.~Schrader, \emph{Axioms for Euclidean Green's functions},
Commun. Math. Phys. \textbf{31}, 83--112 (1973).

\bibitem{OS1975}
K.~Osterwalder and R.~Schrader, \emph{Axioms for Euclidean Green's functions
II}, Commun. Math. Phys. \textbf{42}, 281--305 (1975).

\bibitem{FILS1978}
J.~Fr\"ohlich, R.~Israel, E.~H.~Lieb, and B.~Simon, \emph{Phase transitions
and reflection positivity. I. General theory and long range lattice models},
Commun. Math. Phys. \textbf{62}, 1--34 (1978).

\bibitem{Akhiezer}
N.~I.~Akhiezer, \emph{The Classical Moment Problem}, Oliver \& Boyd (1965).

\bibitem{TitchmarshZeta}
E.~C.~Titchmarsh, revised by D.~R.~Heath-Brown, \emph{The Theory of the Riemann
Zeta-Function}, second edition, Oxford University Press (1986).

\bibitem{Eriksson0064}
L.~Eriksson, \emph{The Heisenberg cut as a resource boundary: an operational
outlook from coherence maintenance costs}, ai.viXra:2512.0064v2 (2026),
\href{https://github.com/lluiseriksson/aqft-split-inclusion-series/tree/main/papers/0064-heisenberg-cut}
{source and verification repository}.

\bibitem{Eriksson0070}
L.~Eriksson, \emph{Stress testing the rate inheritance principle: spectral
decoherence rates and an operational resource horizon}, ai.viXra:2512.0070v2
(2026),
\href{https://github.com/lluiseriksson/aqft-split-inclusion-series/tree/main/papers/0070-rate-inheritance}
{source and verification repository}.

\bibitem{ErikssonProgramme}
L.~Eriksson, \emph{The Eriksson Programme: a Lean 4/Mathlib formalization
toward the Yang--Mills mass gap---with honest goalposts} (2026),
\href{https://github.com/lluiseriksson/THE-ERIKSSON-PROGRAMME}
{GitHub repository}.

\bibitem{PrimeResolvent}
L.~Eriksson, \emph{Riemann prime--resolvent programme} (2026),
\href{https://github.com/lluiseriksson/riemann-prime-resolvent}
{GitHub repository}.

\end{thebibliography}

\end{document}
